\documentclass[troisieme]{fiche}
\metadonnees{9}{Bright River}{Bloc B — Ce qu'on peut, ce qu'on doit}

\begin{document}
\entete

\objectifs{%
\textbf{Langue} : \emph{should}, \emph{had better} ; donner un conseil, proposer.\par
\textbf{Texte} : L.\,M. Montgomery, \emph{Anne of Green Gables} (1908), ch. \textsc{ii}.\par
\textbf{Durée} : 1 h — exercices 20 min, texte et questions 25 min, version 15 min.}

%% ===================================================================
\exercice[12 min]{\emph{Should} et \emph{had better}}

\rappel{\textbf{Rappel.} \emph{Should} conseille. \emph{Had better} avertit : il sous-entend
qu'il y aura des ennuis sinon.}

\consigne{Complétez avec \emph{should}, \emph{shouldn't} ou \emph{had better}.}

\begin{anglais}
\begin{enumerate}[label=\arabic*.,itemsep=2.5mm,leftmargin=*]
  \item You \trou{should} write to your aunt; she is waiting for news.
  \item You \trou{had better} take the bag yourself --- the handle pulls out.
  \item He \trou{shouldn't} leave a child alone at the station.
  \item We \trou{had better} hurry, or we will miss the train.
  \item You look tired; you \trou{should} go to bed early tonight.
  \item She \trou{shouldn't} talk so much to a stranger.
  \item You \trou{had better not} tell her tonight; wait until tomorrow.
  \item I think we \trou{should} ask Mrs Spencer what happened.
  \item Children \trou{shouldn't} be sent away without an explanation.
  \item It is eight miles: you \trou{had better} start now.
\end{enumerate}
\end{anglais}

\note[Deux détails de forme]{La négation de \emph{had better} est \emph{had better not}, et
jamais \emph{had not better}. Ni \emph{should} ni \emph{had better} ne sont suivis de
\emph{to} : \emph{you should go}, \emph{you had better go}.}

%% ===================================================================
\exercice[8 min]{Proposer, conseiller}

\rappel{\textbf{Rappel.} Pour proposer : \emph{let's} + base verbale, \emph{shall we\ldots ?},
\emph{why don't you\ldots ?}, \emph{how about} + \emph{-ing}.}

\consigne{Traduisez en anglais.}

\begin{anglais}
\begin{enumerate}[label=\arabic*.,itemsep=3mm,leftmargin=*]
  \item Allons-y.
    \lignes{1}
    \corr{Let's go.}
  \item Et si on prenait le train ?
    \lignes{1}
    \corr{How about taking the train? \emph{ou} Why don't we take the train?}
  \item Tu devrais lui parler.
    \lignes{1}
    \corr{You should talk to her.}
  \item Pourquoi ne demandes-tu pas à ton oncle ?
    \lignes{1}
    \corr{Why don't you ask your uncle?}
  \item On y va ? \emph{(proposition)}
    \lignes{1}
    \corr{Shall we go?}
  \item Tu ferais mieux de te dépêcher.
    \lignes{1}
    \corr{You had better hurry.}
\end{enumerate}
\end{anglais}

\note[Après \emph{how about}]{Toujours un gérondif : \emph{how about taking the train?} On ne
dit pas \emph{how about to take}. Après \emph{why don't you}, en revanche, la base verbale :
\emph{why don't you ask?}}

%% ===================================================================
\newpage
\section{Texte}

\noindent\small\textit{Matthew Cuthbert et sa sœur Marilla, âgés et sans enfants, ont demandé
à l'orphelinat un garçon pour les aider à la ferme. Matthew, qui est timide, vient le chercher
à la gare de Bright River. Sur le quai l'attend une fillette de onze ans.}\normalsize

\begin{texte}
Matthew, however, was spared the \gl[an ordeal]{ordeal}{une épreuve} of speaking first, for as
soon as she concluded that he was coming to her she stood up, grasping with one thin brown
hand the handle of a \gl[shabby]{shabby}{miteux} old-fashioned
\gl[a carpet-bag]{carpet-bag}{un sac de voyage en tapisserie}; the other she held out to him.

``I suppose you are Mr. Matthew Cuthbert of Green Gables?'' she said in a peculiarly clear,
sweet voice. ``I'm very glad to see you. I was beginning to be afraid you weren't coming for
me and I was imagining all the things that might have happened to prevent you. I had made up
my mind that if you didn't come for me tonight I'd go down the track to that big wild
cherry-tree at the bend, and climb up into it to stay all night. I wouldn't be a bit afraid,
and it would be lovely to sleep in a wild cherry-tree all white with
\gl[bloom]{bloom}{la floraison} in the \gl[moonshine]{moonshine}{le clair de lune}, don't you
think? You could imagine you were \gl[to dwell]{dwelling}{habiter} in marble halls, couldn't
you? And I was quite sure you would come for me in the morning, if you didn't tonight.''

Matthew had taken the \gl[scrawny]{scrawny}{maigrichon} little hand awkwardly in his; then and
there he decided what to do. He could not tell this child with the glowing eyes that there had
been a mistake; he would take her home and let Marilla do that. She couldn't be left at Bright
River anyhow, no matter what mistake had been made, so all questions and explanations might as
well be \gl[to defer]{deferred}{différer, remettre} until he was safely back at Green Gables.

``I'm sorry I was late,'' he said shyly. ``Come along. The horse is over in the yard. Give me
your bag.''

``Oh, I can carry it,'' the child responded cheerfully. ``It isn't heavy. I've got all my
worldly goods in it, but it isn't heavy. And if it isn't carried in just a certain way the
handle pulls out --- so I'd better keep it because I know the exact
\gl[a knack]{knack}{un tour de main} of it. It's an extremely old carpet-bag. Oh, I'm very
glad you've come, even if it would have been nice to sleep in a wild cherry-tree. We've got to
drive a long piece, haven't we? Mrs. Spencer said it was eight miles. I'm glad because I love
driving.''
\end{texte}
\source{L.\,M. Montgomery, \emph{Anne of Green Gables}, Boston, Page, 1908, ch.~\textsc{ii}.}

\partiel[15 min]{Questions}
\consigne{Répondez aux deux premières questions en français, aux deux dernières en anglais.}

\begin{enumerate}[label=\arabic*.,itemsep=2.5mm,leftmargin=*]
  \item Qu'avait prévu la fillette au cas où personne ne serait venu la chercher ?
    \lignes{2}
    \corr{Descendre la voie jusqu'au grand cerisier sauvage du tournant, y grimper et y
    passer la nuit. Elle affirme qu'elle n'aurait pas eu peur du tout et que cela aurait même
    été charmant, sous la floraison blanche et au clair de lune.}
  \item Quelle décision Matthew prend-il, et pourquoi la prend-il si vite ?
    \lignes{3}
    \corr{Emmener l'enfant à Green Gables et laisser Marilla lui annoncer qu'il y a eu une
    erreur. Il ne peut pas le dire lui-même à une petite fille aux yeux brillants, et on ne
    peut pas la laisser seule à la gare : le reste attendra.}
  \item \en{Anne insists three times that the bag is not heavy. What is she really doing?}
    \lignes{3}
    \corr{She is showing that she is no trouble and asks for nothing — and at the same time
    she cannot stop talking. \emph{I've got all my worldly goods in it} says, without
    complaining, that she owns almost nothing.}
  \item \en{Find in the text one sentence with \emph{had better} and one with \emph{could}.
    Who says each of them?}
    \lignes{2}
    \corr{\emph{So I'd better keep it because I know the exact knack of it} : c'est Anne, qui
    se donne un conseil à elle-même. \emph{He could not tell this child\ldots} : c'est le
    narrateur, sur Matthew — une impossibilité, non une interdiction.}
\end{enumerate}

\note[Ce que le texte ne dit pas]{Matthew ne dit rien de l'erreur, et Anne ne se doute de
rien. Le lecteur sait ce que l'enfant ignore : c'est de cet écart que vient l'émotion du
chapitre, et non de ce qui est dit.}

%% ===================================================================
\section{Version}
\consigne{Traduisez le troisième paragraphe du texte, de \emph{Matthew had taken} à
\emph{back at Green Gables}.}

\lignes{12}

\corr{\textbf{Proposition de traduction.} Matthew avait pris gauchement dans la sienne la
petite main maigrichonne ; sur-le-champ, il décida de ce qu'il ferait. Il ne pouvait pas dire
à cette enfant aux yeux brillants qu'il y avait eu une erreur ; il l'emmènerait chez lui et
laisserait Marilla s'en charger. On ne pouvait de toute façon pas l'abandonner à Bright River,
quelle que fût l'erreur commise ; aussi toutes les questions et toutes les explications
pouvaient-elles aussi bien attendre qu'il fût rentré sain et sauf à Green Gables.}

\note[Choix de traduction 1]{\emph{he would take her home and let Marilla do that} :
\emph{would} est ici le futur vu du passé, donc un conditionnel en français — \og il
l'emmènerait \fg. Ce n'est pas une habitude ni une politesse.}

\note[Choix de traduction 2]{\emph{She couldn't be left at Bright River} : passif sans
complément d'agent. Le français passe par \og on \fg{} : \og on ne pouvait pas
l'abandonner \fg.}

\note[Choix de traduction 3]{\emph{might as well be deferred} : \emph{might as well} =
\og autant \fg, \og aussi bien \fg. Locution à reconnaître d'un bloc ; traduire \emph{might}
seul donnerait un contresens.}

%% ===================================================================
\partiel{Lexique à mémoriser}
\begin{multicols}{2}\small
\begin{itemize}[label=--,itemsep=0.8mm,leftmargin=*]
  \item \textbf{shy} — timide
  \item \textbf{awkwardly} — gauchement
  \item \textbf{a handle} — une poignée
  \item \textbf{to hold out} — tendre (la main)
  \item \textbf{to make up one's mind} — se décider
  \item \textbf{to prevent} — empêcher
  \item \textbf{a mistake} — une erreur
  \item \textbf{glowing} — brillant, rayonnant
  \item \textbf{a yard} — une cour
  \item \textbf{heavy} — lourd
  \item \textbf{goods} — les biens, les marchandises
  \item \textbf{to drive} — conduire ; drove, driven
\end{itemize}
\end{multicols}

\end{document}
